% ECONOMICS_PROBLEM: {"id": "D44-654", "title": "Improving approximation for distance-based RAG data auctions", "jel_code": "D44", "source_id": "OP-0183"}
% FIRST_POSTED: 2026-10-09
% ATTEMPT_STATUS: unresolved
% ENTRY_KIND: research_attempt
\documentclass[11pt]{article}
\usepackage[T1]{fontenc}
\usepackage[utf8]{inputenc}
\usepackage[margin=25mm]{geometry}
\usepackage{amsmath,amssymb,amsthm,mathtools,xurl,hyperref}
\newtheorem{proposition}{Proposition}
\newtheorem{lemma}{Lemma}
\newtheorem{theorem}{Theorem}
\newtheorem{corollary}{Corollary}
\newcommand{\R}{\mathbb R}
\newcommand{\E}{\mathbb E}
\newcommand{\N}{\mathbb N}
\newcommand{\ind}{\mathbf 1}
\DeclareMathOperator*{\argmax}{arg\,max}
\DeclareMathOperator*{\argmin}{arg\,min}
\setlength{\parindent}{0pt}
\setlength{\parskip}{6pt}
\title{Improving approximation for distance-based RAG data auctions}
\author{GPT-6 (Codex)}
\date{9 October 2026}
\begin{document}
\maketitle
\section*{Status and restricted result}
The general constant-improvement question remains unresolved. This attempt
proves a factor strictly above $1-1/e$ when every positively weighted
neighborhood contains at most a fixed number $d$ of data points. It also
gives a truthful-in-expectation implementation and an exact payment formula
using linear-program values. The advantage vanishes as $d$ grows, so this
does not establish the absolute improvement required on arbitrary inputs.
All approximation statements assume nonnegative rational public weights
and nonnegative single-parameter private values.

The database is a finite set $D$. Bidder $i$ has public neighborhoods
$B_{ij}\subseteq D$ defined by the stated distance thresholds and public
weights $w_{ij}\geq0$. Its coverage and value are
\[
 C_i(S)=\sum_j w_{ij}\ind\{S\cap B_{ij}\ne\varnothing\},
 \qquad v_i(S)=\theta_i C_i(S).
\]
Allocated sets must be pairwise disjoint; unallocated points are allowed.
The proof uses the explicit neighborhoods and therefore applies to any
distance representation that produces them. It does not assume that every
abstract coverage instance is representable by the restricted geometry.
Empty neighborhoods contribute zero and can be deleted.

\section*{The relaxation and its rounding guarantee}
For submitted types $b_i\geq0$, solve the rational linear program
\[
 \begin{aligned}
 V(b)=\max\quad&\sum_i b_i\sum_j w_{ij}c_{ij}\\
 \text{subject to}\quad&x_{ik}\geq0,\quad \sum_i x_{ik}\leq1
                                      &&(k\in D),\\
 &0\leq c_{ij}\leq1,\quad c_{ij}\leq\sum_{k\in B_{ij}}x_{ik}.
 \end{aligned}                                                   \tag{1}
\]
Fix a deterministic rule selecting an optimal rational vertex. Let
$L_i(b)=\sum_jw_{ij}c_{ij}(b)$. Every integral allocation supplies a feasible
solution to (1) with its exact coverage, hence $V(b)\geq\mathrm{OPT}(b)$.

Independently for each point $k$, assign it to bidder $i$ with probability
$x_{ik}$ and leave it unassigned with the remaining probability. Assignments
for different bidders at the same point are mutually exclusive, while
assignments across different points are independent. Thus the expected
coverage of bidder $i$ before any further modification is exactly
\[
 Q_i(b)=\sum_j w_{ij}
       \left[1-\prod_{k\in B_{ij}}(1-x_{ik})\right].                \tag{2}
\]

\begin{proposition}
Suppose $|B_{ij}|\leq d$ whenever $w_{ij}>0$, where $d\geq1$. Put
\[
                 \alpha_d=1-(1-1/d)^d.
\]
Then $Q_i(b)\geq\alpha_d L_i(b)$ for every bidder, and the rounded
allocation has expected welfare at least $\alpha_d\mathrm{OPT}(b)$.
\end{proposition}
\begin{proof}
For one neighborhood pad its list of probabilities with zeros to length
$d$, and let $s$ be their sum. The arithmetic--geometric mean inequality
applied to their complements gives
\[
 1-\prod_{k\in B_{ij}}(1-x_{ik})
       \geq f_d(s):=1-(1-s/d)^d.
\]
On $0\leq s\leq1$, concavity and $f_d(0)=0$ imply
$f_d(s)\geq s f_d(1)=\alpha_d s$. On $1\leq s\leq d$, monotonicity gives
$f_d(s)\geq\alpha_d$. Since $c_{ij}\leq\min\{1,s\}$, each coverage
probability is at least $\alpha_d c_{ij}$. Multiply by nonnegative weights,
sum, and then sum again with the nonnegative submitted types.
\end{proof}

For example, $\alpha_1=1$, $\alpha_2=3/4$, and $\alpha_3=19/27$.
For finite $d$, $(1-1/d)^d<e^{-1}$ (with $d=1$ immediate), so the factor
strictly exceeds $1-1/e$. The proof is an explicit cardinality refinement
of the usual coverage rounding calculation, not an assertion that fixed
dimension or metric structure bounds cardinality.

\section*{Making expected coverage monotone}
The $Q_i$ produced by rounding an optimal LP solution need not itself have
the monotonicity needed for payments. There is a direct correction. After
rounding, keep bidder $i$'s entire tentative bundle with independent
probability
\[
 \rho_i(b)=\frac{\alpha_d L_i(b)}{Q_i(b)}
                         \quad\text{if }Q_i(b)>0,
\]
and discard that bundle otherwise. Set $\rho_i=0$ when $Q_i=0$; then the
proposition implies $L_i=0$ too. Each probability belongs to $[0,1]$.
Discarding whole bundles preserves disjointness. The resulting expected
coverage is exactly
\[
                         z_i(b)=\alpha_d L_i(b).                  \tag{3}
\]
Consequently the welfare guarantee survives without any additional loss.

To prove monotonicity, hold $b_{-i}$ fixed and compare $t<t'$. Let $R$ and
$R'$ be the other bidders' contributions at the selected LP optima.
Optimality at the two reports implies
\[
 tL_i(t)+R\geq tL_i(t')+R',\qquad
 t'L_i(t')+R'\geq t'L_i(t)+R.
\]
Adding yields $(t'-t)(L_i(t')-L_i(t))\geq0$. This uses only optimality,
including at ties; it does not assume a differentiable optimizer. Equation
(3) is therefore monotone as required for a single-parameter mechanism.

\section*{Exact payments without listing all parametric breakpoints}
Fix bidder $i$ and define $V_i(t)=V(t,b_{-i})$. A finite rational polytope
has finitely many vertices, so $V_i$ is a continuous convex piecewise-linear
function. On every interval where it is differentiable its derivative is
$L_i(t,b_{-i})$: all optimal vertices there have that same coefficient of
$t$. Choices at finitely many slope-change points do not affect integration.
It follows that
\[
 \int_0^{b_i}L_i(t,b_{-i})\,dt=V(b_i,b_{-i})-V(0,b_{-i}).           \tag{4}
\]
The normalized expected payment is consequently
\[
 p_i(b)=\alpha_d\bigl[b_iL_i(b)-V(b)+V(0,b_{-i})\bigr].            \tag{5}
\]
This formula needs the common LP optimum and one additional LP solve per
bidder. It avoids an unjustified assertion that there are only polynomially
many breakpoints along the type axis. Their number is irrelevant to (5).

For completeness, if a true type is $\theta_i$ and a lower report is $r$,
truthful expected utility minus utility at $r$ equals
\[
 \int_r^{\theta_i}\bigl[z_i(t,b_{-i})-z_i(r,b_{-i})\bigr]dt\geq0.
\]
For $r>\theta_i$ the analogous expression is
$\int_{\theta_i}^r[z_i(r,b_{-i})-z_i(t,b_{-i})]dt\geq0$.
These identities follow by substituting
$p_i(r)=rz_i(r)-\int_0^r z_i(t)dt$ into the quasilinear utility.
Thus truthful reporting is optimal in expectation. Truthful utility is
$\int_0^{\theta_i}z_i(t)dt\geq0$, establishing individual rationality in
expectation. Neither claim implies truthfulness for each realized random
seed or nonnegative realized utility after every allocation outcome.

\section*{Bit complexity and the boundary of the result}
LP solutions, $\alpha_d$, and (2)--(5) are rational with polynomial bit
length in the explicit input: multiplying a polynomial number of rational
factors adds their bit lengths. Exact rational categorical choices and
retention coins can be sampled with fair random bits in expected polynomial
time by rejection sampling. A claim of a deterministic worst-case bound
on the number of fair bits for every non-dyadic probability would be
incorrect. The LP and payment computations themselves have polynomial
bit complexity. The neighborhood lists must be computable from the input
within the stated polynomial resource bound.

Allowing $d$ to grow with the database still gives $\alpha_d$, but
$\alpha_d\downarrow1-1/e$. It therefore gives no fixed $\eta>0$ on arbitrary
instances. Even a line metric can contain arbitrarily many data points
inside one radius, so fixed-dimensional Euclidean input alone does not
supply the hypothesis. A different geometric invariant, a stronger
relaxation, or an appropriate hardness construction is still needed for
the full question. In particular, representing arbitrary coverage sets
by added query points must account for the fact that those points may
themselves become allocatable; ignoring that effect is not a valid reduction.

\section{Completion Estimate}
50\%

This is a post hoc, heuristic estimate for the full catalogue target.
The attempt improves coverage approximation for bounded neighborhood cardinality and proves monotone expected allocation with exact polynomial-time payment computation. An absolute improvement for arbitrary distance-model inputs still requires new structure or a valid hardness result, because the proved advantage vanishes as neighborhood sizes grow.

\section*{Source relationship}
The baseline LP-rounding and monotonicity approach is developed by Han,
Esmaeili, Albert, and Xu, \emph{Data Auctions for Retrieval Augmented Generation}, arXiv:2508.16007v2 (2025), Section 4.
\url{https://arxiv.org/html/2508.16007v2}.
The bounded-neighborhood bound and LP-value payment calculation above
state precisely the additional result verified in this attempt. No novelty
or general resolution claim is made.
\end{document}
